\subsection{Compact $E_8$ compression and the role of $\nu^c$}
\begin{equation}
  \boxed{\ \begin{aligned}
    |R(E_8)|&=\dim S^+\cdot\gcar\cdot\Nfam=16\cdot5\cdot3=240,\\
    \dim E_8&=\dim S^+(\AH+\AL)+(\gcar+\Nfam)=240+8=248
  \end{aligned}\ }
\end{equation}
with $\Nfam=(2^{\gcar-1}-1)/\gcar=15/5=3$ and $\operatorname{rank}E_8=\gcar+\Nfam=8$. \mE. The
right-handed neutrino $\nu^c$ is the Pascal singleton $\Lambda^0E=(1,1)_0$ that closes $16=1+5+10$;
removing it breaks the truncation and shifts $\alpha^{-1}$ by $2.16\times10^{-3}$.

% =====================================================================
\section{The $D_5\times A_3$ compiler: $E_8$ as a closure, not an input}
\label{sec:d5a3}
% =====================================================================
This section upgrades the $E_8$ atlas (Section~\ref{sec:e8}) from a counting coincidence to an explicit
lattice construction. The whole content is consolidated and verified in the standalone companion note
\emph{``The Even Carrier Code and the $\mathbb Z_{30}$ Coxeter Compiler.''} The thesis: the carrier
code is the $D_5$ half-spinor, the family geometry is an $A_3$, and $E_8$ is the standard unimodular
glue closure of $D_5\oplus A_3$ --- with the four punctures $\mu_4$ as the glue code.

\TFPTeightglue

\par\smallskip\noindent\textbf{The compiler in one line:}
\begin{keyeq}
$\displaystyle \text{five carrier slots}\Rightarrow D_5,\quad
  \text{four family corners }\mu_4\Rightarrow A_3,\quad
  (D_5\oplus A_3)+\mu_4\text{-glue}\Rightarrow E_8.$
\end{keyeq}

\begin{figure}[ht]
\centering
\begin{tikzpicture}[>=Stealth]
  \node[cbox] (seam) at (0,2)    {boundary\\(seam)};
  \node[cbox] (c3)   at (3.4,2)  {$c_3=\tfrac1{8\pi}$};
  \node[cbox] (cont) at (8.6,2)  {continuous readouts:\\$\alpha^{-1},\ \lambda_C,\ \rho_\Lambda,\ H_0,\ n_s,\dots$};
  \node[dbox] (slots) at (0,0.4)  {five carrier slots\\$C^+$ (even code, $16$)};
  \node[dbox] (d5)    at (3.4,0.4){$D_5=\mathfrak{so}_{10}$\\half-spinor; $W(D_5){=}1920$};
  \node[abox] (corn)  at (0,-1.3) {four corners $\mu_4$\\$\mathbb P^1\setminus\mu_4$};
  \node[abox] (a3)    at (3.4,-1.3){$A_3=\mathfrak{su}_4$\\$12$ roots, rank $3{=}N_{\mathrm{fam}}$};
  \node[draw=purple!55!black, fill=purple!6, rounded corners, align=center,
        font=\scriptsize, inner sep=3pt] (glue) at (6.1,-0.45) {$+\,\mu_4$ glue\\$\mathbb Z_4$};
  \node[ebox] (e8) at (8.6,-0.45) {$E_8$\\$240$ roots, $\det 1$};
  \draw[flow] (seam)--(c3);   \draw[flow] (c3)--(cont);
  \draw[flow] (slots)--(d5);  \draw[flow] (corn)--(a3);
  \draw[flow] (d5)--(glue);   \draw[flow] (a3)--(glue);  \draw[flow] (glue)--(e8);
\end{tikzpicture}
\caption{The explanation hierarchy. The boundary fixes the seam seed $c_3$; the five-slot even
carrier code is the $D_5$ half-spinor; the four punctures $\mu_4$ give $A_3$; their common
$\mathbb Z_4$ discriminant ($\mu_4$) glues $D_5\oplus A_3$ into the unimodular $E_8$. $E_8$ is the
\emph{closure}, not the input; $c_3$ dresses the continuous observables.}
\label{fig:pipeline}
\end{figure}

\par\smallskip\noindent\textbf{The glue is \emph{one} operation in two categories.} The same
order-$|\mu_4|{=}4$ discriminant glue closes $D_5\oplus A_3$ into $E_8$ \emph{twice over}: once on
\emph{root lattices} (the unimodular pushout $D_5\oplus A_3+\mu_4\Rightarrow E_8$, \vref{v1}) and once on
\emph{chiral conformal nets} (the simple-current / anyon-condensation extension
$(D_5)_1\otimes(A_3)_1\rtimes\mu_4\cong(E_8)_1$, \vref{v92}/\vref{v281}). These are not two coincidences but
the two functorial faces of a single commuting square (\vref{v234}/\vref{v277}/\vref{v260}); the boundary
realisation on which the strict-TOE residual turns (\texttt{SEAM.EQUIV.01}) is exactly the statement that the
raw seam realises this square --- now \emph{closed modulo cited theorems on its MMST route}
\texttt{SEAM.EQUIV.MMST.01} (the lattice/$S_3$ stack pins it at
every computable level, Lean \texttt{FORM.SEAM.MMST.01}; only the cited MMST/OS continuum-existence theorems stay
\mO, \vref{v367}/\vref{v368}/\vref{v376}--\vref{v379}/\vref{v336}; the twistor route
\texttt{SEAM.EQUIV.TWISTOR.01} and the unconditional parent stay \mO). The extension leg of that citation package is
itself re-founded at net level (\veri{v469\_seam\_crossedproduct\_route.py}): the $128$-spinor step is the
\emph{local} $\Z_2$ simple-current crossed product ($h_s(SO(16)_1)=16/16=1\in\Z$, the Longo--Rehren locality
criterion; KLM $\mu=4/2^2=1$), on 1995--2001 peer-reviewed subfactor theory, with the lattice-VOA preprint route
demoted to an independent second witness, and the realisation input reduced to invariant level (R1$'$).
Premise (i) of that package ($c_-{=}8$) now also carries a \emph{ground-state} witness
(\veri{v489\_seam\_modular\_commutator.py}): the Kim--Shehab--Kim modular commutator
$J(A,B,C)=(\pi/3)c_-$, evaluated on the exact BdG ground state of the collar model, gives $c_-=8=\gcar+\Nfam$
for the $16$-layer carrier (per layer $0.500000$ at $L{=}48$; trivial and orientation-flip controls pass) ---
independent of both the Chern route (\vref{v367}) and the KLM count (\vref{v378}).
\texttt{SEAM.EQUIV.01} stays \mO.
\emph{Follow-up (2026-08-03, the $G_{\mathrm{net}}$ witnesses):} the
index-$4$ statement now carries two \emph{exact} arithmetic/index
witnesses.  The NS/R grading \emph{is} the parity character of
$E_8(\Z[i])/(1{+}i)=\mathbb F_2^4$ --- $(1{+}i)L$ lies in the even
sublattice (proved on a $\Z$-basis), all $15\times16$ root classes
are parity-pure ($7$ NS $+$ $8$ R), so the Ramond projection is
$(1{+}i)$-\emph{adic} at the one ramified edge of the
$\Z[i]$-E8 Hecke tower (\veri{v722\_phys\_ramified\_ns\_r.py},
$9/9$, \textsc{ramified-sees-nsr}); and on the CAR ladder the
$\mu_4$-average expectation has optimal Pimsner--Popa constant
exactly $1/4$ and Watatani index exactly $4$ via an explicit
quasi-basis with \emph{forced} weights, with $\mu_4$ \emph{derived}
from the clock and the Ramond sector healed state-preservingly by
the Klein zero-mode dressing (\veri{v726\_phys\_car\_pp\_index.py},
$57/57$, \textsc{t2-slice-go}).  The remaining half is the
\emph{identification} of the ladder Q-system with the constructed
$\mathbb C[\Z_4]$ object (\texttt{GNET.RAMIFIED.01}, kills
registered); \texttt{GATE.METRIC.08}/\texttt{.10} typing unchanged.
\emph{Follow-up (2026-08-05, the continuum strand):} the
finite-level precursors of the GNS limit are measured
\emph{coherent} (\veri{v779\_gnet\_gns\_arf.py}, $27/27+27/27$):
the Haar scaling tower intertwines exactly, the index-4
expectation squares commute exactly (the finite half of the
Q-system identification), the state converges geometrically, and
the exact dim law $d_c(\ell)=4^{\ell-1}+2^{\ell-1}\cos(\pi c/2)$
types the $\mathbb C[\Z_4]$ obstruction at exactly $2^{-\ell}$
(explaining the $\ell=1$ three-sector anomaly); the honest
negative: the $16=1{+}5{+}10$ counting echo with the Arf geometry
is exact but \emph{non-geometric}, and the two fours (index $4$
vs channel $4/49$) stay in their registers.  The round-21
martingale probe measures the limit-existence hypotheses of the
continuum strand as met
(\veri{v802\_gnet\_martingale.py}, $15/15$,
\textsc{martingale-hypotheses-met}; the analytic remainder stays
typed).  Typing unchanged.

\TFPTdoublepushout

\subsection{Carrier $=D_5=\mathfrak{so}_{10}$ (proved)}
The $16$ even-occupation states $C^+=\{x\in\mathbb F_2^{\gcar}:|x|\ \text{even}\}$ are exactly the
weights of the $D_5$ half-spinor $\tfrac12(\pm1,\dots,\pm1)$ (even number of minus signs). The
flip-weight-$4$ relation is the \textbf{Clebsch graph} $\mathrm{SRG}(16,5,0,2)$ and the flip-weight-$2$
relation is its complement $\mathrm{SRG}(16,10,6,6)$; the \emph{affine} automorphism group of the scheme
(translations $\rtimes$ coordinate permutations) is the $D_5$ Weyl group,
\begin{equation}
  \boxed{\ \operatorname{Aut}_{\mathrm{aff}}(C^+)=C^+\rtimes S_5,\quad |\operatorname{Aut}_{\mathrm{aff}}(C^+)|=2^{\gcar-1}\gcar!=1920=|W(D_5)|=r\,|R(E_8)|\ }\qquad\mE ,
\end{equation}
(the linear code automorphism group alone is only $S_5$; the order-$1920$ group is the affine/scheme one)
and the Clebsch edge count is the flavor budget $|E(R_4)|=\tfrac{16\cdot5}{2}=40=|R(D_5)|=\sum_{f,j}L_{f,j}$.
\emph{Round-22 grading (2026-08-06):} the Arf split $16=1+\bar5+10$ \emph{is} the cut
classification $(0|5)+(1|4)+(2|3)$ of $K_5$ (the cut map injective-isometric on all $256$ cells),
and the Clebsch edges grade as $40=5+20+15$ over the Arf shells with the within-$10$ block equal to
the Petersen graph --- an Arf/cut geometry, exact but carrying no matter semantics
(\veri{v810\_cutcode\_moment.py}) \mE.

\subsection{Family $=A_3=\mathfrak{su}_4$ from $\mathbb P^1\setminus\mu_4$ (proved, with the metric settled)}
The carrier geometry fixes $X_f^\circ\cong\mathbb P^1\setminus\mu_4$, $H_1=\mathbb Z^3$, rigid $D_4$ generated by
$z\mapsto iz,\ z\mapsto1/z$. The twelve puncture-difference cycles $\gamma_i-\gamma_j$ carry the
residue pairing $\operatorname{Res}(\omega_i,\gamma_j)=\delta_{ij}$ and form the $A_3$ root system:
\begin{equation}
  \boxed{\ \#\{\gamma_i-\gamma_j\}=12=|R(A_3)|=\dimg,\quad \operatorname{rank}=3=\Nfam,\quad
  \mathrm{Gram}=\mathrm{Cartan}(A_3),\ \det=4\ }\qquad\mE .
\end{equation}
The corner rotation $z\mapsto iz$ acts on $H^1=\mathbb C^3$ as the $A_3$ Coxeter element (order
$4=h(A_3)=|\mu_4|$, spectrum $\{i,-1,-i\}$, $\chi=\Phi_4\Phi_2$), an isometry of the Cartan form ---
the exact analogue of the $E_8$ Coxeter element reading $\Phi_{30}$.
\par\smallskip\noindent\textbf{The $\mu_4$ clock is the order-$4$ character of the $E_8$ cycle, and $(2,3,5)$
is the icosahedron (\veri{v223\_coxeter\_totative\_clock.py}, \veri{v219\_icosahedral\_mckay.py}).} The two
``Coxeter'' clocks are one object read at two scales. The $E_8$ exponents are the totatives
$(\Z/30)^\times=\{1,7,11,13,17,19,23,29\}$, and the multiplicative order-$4$ element $7$ ($7^2{=}19$,
$7^4{=}1\bmod30$) generates a literal $\mu_4$ clock $\langle7\rangle=\{1,7,13,19\}$ that permutes the four
conjugate Coxeter $2$-planes --- the same $\mu_4$ as the seam corner rotation. And the period
$30=2{\cdot}3{\cdot}5$ itself is not a free factorisation: by the McKay correspondence $E_8$ is the
exceptional \emph{top} of the tower of finite $SU(2)$ subgroups ($2T{\to}\hat E_6$, $2O{\to}\hat E_7$,
$2I{\to}\hat E_8$), so its branch data is the \emph{icosahedron} and the three atoms $(|\Z_2|,\Nfam,\gcar)
=(2,3,5)$ are its rotation-axis orders. The binary icosahedral group $2I$ (order $120=|R^+(E_8)|$) has nine
irreducible-representation degrees $\{1,2,2,3,3,4,4,5,6\}$ equal to the affine-$E_8$ Kac marks, with
$\sum d_i=30=h(E_8)$ and $\sum d_i^2=120$. This is a \emph{backward certificate} of the closed $E_8$ \mE\
(it explains the $(2,3,5)$ choice \mC; it does \emph{not} replace the P2 input).
\par\smallskip\noindent\textbf{The metric, settled (former \mC).}
The transcendental geometric (Green) metric $G_{ij}=-\log|p_i-p_j|$ on $\mu_4$ is only
$D_4$-symmetric, with a regularisation-independent anisotropy $-\log2$ on the root space, so it is
\emph{not} proportional to the $A_3$ Cartan form. But it does not need to be: the $E_8$-relevant $A_3$
is the canonical \emph{residue} pairing (integral, $S_4$-symmetric, $=\mathrm{Cartan}(A_3)$), and the
geometric metric is a $D_4$-equivariant deformation preserving the root system, Weyl group and Coxeter
element. So the lattice structure is exact \mE; the analytic metric is a characterised deformation,
not an open gap.
\par\smallskip

\subsection{The $\mu_4$ glue and the $5/4+3/4=2$ discriminant theorem (proved)}
The two factors have equal discriminant groups and $E_8$ is unimodular, so the glue index is exactly
$|\mu_4|$:
\begin{equation}
  \operatorname{disc}(D_5)=\operatorname{disc}(A_3)=\mathbb Z_4,\qquad
  [E_8:D_5\oplus A_3]=\sqrt{\tfrac{\det D_5\det A_3}{\det E_8}}=\sqrt{\tfrac{4\cdot4}{1}}=4=|\mu_4| .
\end{equation}
The discriminant-form norms sum to the $E_8$ root norm, and --- strikingly --- they are pre-existing
TFPT constants:
\begin{equation}
  \boxed{\ \begin{aligned}
    &q(D_5)=\tfrac54=\frac{\dtwo}{\dtop^2},\qquad q(A_3)=\tfrac34=\xi_{\mathrm{tree}},\\
    &q(D_5)+q(A_3)=2=|\text{$E_8$ root}|^2
  \end{aligned}\ }\qquad\mE\ \veri{v1\_e8\_glue.py}
\end{equation}
This is literally the $E_8$ spinor-root split $\tfrac12(\pm)^8=\tfrac12(\pm)^5\oplus\tfrac12(\pm)^3$,
$5\cdot\tfrac14+3\cdot\tfrac14=2$.

\begin{figure}[ht]
\centering
\begin{tikzpicture}[>=Stealth]
  \def\sc{3.3}
  \fill[blue!16]   (0,0) rectangle ({1.25*\sc},0.62);
  \fill[orange!20] ({1.25*\sc},0) rectangle ({2*\sc},0.62);
  \draw[line width=0.7pt] (0,0) rectangle ({2*\sc},0.62);
  \draw[line width=0.7pt] ({1.25*\sc},0)--({1.25*\sc},0.62);
  \node at ({0.625*\sc},0.31) {$q_{D_5}=\tfrac54$};
  \node at ({1.625*\sc},0.31) {$q_{A_3}=\tfrac34$};
  \node[below, font=\scriptsize, text=blue!45!black]   at ({0.625*\sc},-0.05) {$=\delta_2/\delta_{\mathrm{top}}^2$ ($D_5$ spinor)};
  \node[below, font=\scriptsize, text=orange!55!black] at ({1.625*\sc},-0.05) {$=\xi_{\mathrm{tree}}$ ($A_3$ corner)};
  \draw[{Stealth}-{Stealth}, line width=0.6pt] (0,1.0)--({2*\sc},1.0);
  \node[above, font=\small] at ({\sc},1.0) {$q_{D_5}+q_{A_3}=2=|E_8\ \text{root}|^2$};
\end{tikzpicture}
\caption{The discriminant-norm glue. The $E_8$ spinor root $\tfrac12(\pm)^8$ (norm $2$) splits across
$D_5\times A_3$ as $5\cdot\tfrac14+3\cdot\tfrac14$; the two halves are the pre-existing TFPT constants
$\delta_2/\delta_{\mathrm{top}}^2=\tfrac54$ and $\xi_{\mathrm{tree}}=\tfrac34$.}
\label{fig:gluenorm}
\end{figure}

\noindent The standard branching is realised explicitly:
\begin{equation}
  \mathfrak e_8=(45,1)\oplus(1,15)\oplus(10,6)\oplus(16,4)\oplus(\overline{16},\overline4),\qquad
  248=45{+}15{+}60{+}64{+}64,\quad 240=40{+}12{+}60{+}64{+}64 ,
\end{equation}

\begin{figure}[ht]
\centering
\begin{tikzpicture}[>=Stealth]
  \node[dbox] (b1) at (0,0)    {$(45,1)$\\$D_5$ adjoint\\norm $2$};
  \node[abox] (b2) at (2.5,0)  {$(1,15)$\\$A_3$ adjoint\\norm $2$};
  \node[cbox] (b3) at (5.0,0)  {$(10,6)=60$\\vector\,$\otimes$\,vector\\$1{+}1=2$};
  \node[ebox] (b4) at (8.1,0)  {$(16,4)\oplus(\overline{16},\overline4)=128$\\spinor\,$\otimes\,\mu_4$\\$\tfrac54{+}\tfrac34=2$};
  \node[font=\small] at (4.0,-1.45)
    {$\dim E_8=248=\underbrace{45{+}15{+}60}_{\mathfrak{so}_{16}\ \mathrm{adjoint}=120}+\underbrace{64{+}64}_{128\ \mathrm{spinor}}$};
\end{tikzpicture}
\caption{The $E_8\supset D_5\times A_3$ branching as norm-$2$ blocks. The adjoint half is the
hypercharge moment $\operatorname{Tr}_{S^+}X^2=120=\dim\mathfrak{so}_{16}$; the spinor half is two
sheets over the four corners $\mu_4$.}
\label{fig:e8blocks}
\end{figure}
and assembling $D_5\oplus D_3\subset D_8$ with the even spinor coset produces all $240$ roots
(norm $2$, Cartan $\det=1$, every root an integer combination of one simple-root basis): \emph{the
$240$-metrisation is no longer blind --- $E_8$ is built}. \mE

\subsection{Glue-norm arithmetic, and three readings of $41$, $60$, $120$}
The two glue norms generate the carrier integers through sum, difference and product. With
$\dim S^+=16$:
\begin{equation}
  \boxed{\ \begin{aligned}
    &16(q_D{+}q_A)=32=2^{\gcar}=h(D_5)h(A_3),\qquad 16(q_D{-}q_A)=8=r(E_8)=h(D_5),\\
    &16\,q_Dq_A=15=\dim S^+{-}1
  \end{aligned}\ }\qquad\mE ,
\end{equation}
so the glue already contains the full carrier exhaustion $2^{\gcar}$, the $E_8$ rank, and the
non-trivial transition count $\AH{+}\AL=15$; and $|\mu_4|=(q_D-1)^{-1}=(1-q_A)^{-1}=4$ is the inverse
displacement of the two norms from the unit norm.

\begin{keybox}[Spine Quotient Lemma: the recurring factors are adjacent quotients of $(2,3,4,5)$ \mE\ (\veri{v74\_compiler\_micro\_lemmas.py})]
The integers of the spine are $(|\Z_2|,\Nfam,|\mu_4|,\gcar)=(2,3,4,5)$. Every ``mysterious'' $O(1)$ factor in
TFPT is an \emph{adjacent quotient} of this chain, not a separate special number:
\[
  \tfrac{|\Z_2|}{\Nfam}=\tfrac23,\quad \tfrac{\Nfam}{|\mu_4|}=\tfrac34=q(A_3),\quad
  \tfrac{|\mu_4|}{\gcar}=\tfrac45,\quad \tfrac{\gcar}{|\mu_4|}=\tfrac54=q(D_5),\quad \tfrac{|\mu_4|}{\Nfam}=\tfrac43 .
\]
These are exactly the factors that recur across the theory: $\lambda_2=(\tfrac23)^6$ (the transport gap
$=$ Page recovery rate, and the Koide target $Q_\star=\tfrac23$); $\varepsilon=\tfrac34\varphi_0=q(A_3)\varphi_0$
(the solar misalignment); $\varphi_0=\tfrac43 c_3+48c_3^4$ with leading $\tfrac43 c_3=\tfrac1{6\pi}$; and
$q(D_5){+}q(A_3)=\tfrac54{+}\tfrac34=2$ (the $E_8$ root norm). So $(\tfrac23,\tfrac34,\tfrac45,\tfrac54,\tfrac43)$
are five projections of one discrete chain --- not scattered constants. The spine is moreover \emph{one}
finite object: as a tetrahedron on $T=\{e_3(a),p_0(a),e_1(a),e_2(a)\}$ its edge and face products and its
volume $120=|R^+(E_8)|=5!$ reproduce the central integer grammar, with the honest negative control that
the graph reading $240=|\mu_4|\cdot|E(K_4)|\cdot|E(K_5)|$ breaks at $K_6$
(\veri{v91\_spine\_tetrahedron.py}).
\end{keybox} The carrier adjoint $\dim D_5=|R(D_5)|+\gcar=45$
then gives a second reading of the $U(1)$ integer,
\begin{equation}
  \boxed{\ 10b_1=\dim D_5-|\mu_4|=45-4=41=\underbrace{40}_{\sum L}+\underbrace{(\gcar-|\mu_4|)}_{=N_\Phi=1}\ }\qquad\mE ,
\end{equation}
i.e.\ the single Higgs doublet is the carrier rank minus the glue corners, $5-4=1$.
\par\smallskip
\noindent\emph{A fourth, metric reading of $41$ (\veri{v222\_cm\_norm\_duality.py}).} On the square seam
($j{=}1728$, complex multiplication by the Gaussian integers $\Z[i]$) the $U(1)$ index is the Gaussian
\emph{norm} of the carrier-glue vector,
\[
  10b_1=N_{\Z[i]}(\gcar{+}i|\mu_4|)=\gcar^2+|\mu_4|^2=5^2+4^2=41 ,
\]
and the hexagonal partner ($j{=}0$, Eisenstein $\Z[\omega]$) reads the same $(3,2)$ carrier split as the
scalaron exponent $N_{\Z[\omega]}(3{+}2\omega)=3^2{-}3{\cdot}2{+}2^2=7$. So $41$ (EM) and $7$ (scalaron)
are the two exceptional CM norms of the one carrier, square vs.\ hexagonal --- not two unrelated $O(1)$
numbers \mE. The start defect
$60$ and the hypercharge moment $120$ are ``bilingual'':
\begin{equation}
\begin{gathered}
  \Dstart=\dim D_5+\dim A_3=45+15=60=\AL\cdot|R^+(A_3)|=10\cdot6,\\[2pt]
  \operatorname{Tr}_{S^+}X^2=120=\underbrace{60}_{\text{adjoint}}+\underbrace{60}_{\text{mixed }(10,6)} ,
\end{gathered}
\end{equation}
and $\dim E_8=248=\operatorname{Tr}_{S^+}X^2+2\dim S^+|\mu_4|=120+2\cdot16\cdot4$. Finally the
hypercharge norm and the EW exponent are $A_3$-rational,
\begin{equation}
  \boxed{\ \gamma_{\mathrm{car}}=\frac{\AL}{|R(A_3)|}=\frac{10}{12}=\frac56,\qquad
  I_1^{\mathrm{EW}}=\frac{|E(K_5)|}{|R(A_3)|}+\frac{1}{|\mu_4||R(A_3)|}=\frac{10}{12}+\frac1{48}=\frac{41}{48}\ }\qquad\mE .
\end{equation}

\noindent\emph{Transport reduction.} The flavor matrix $L$ has its mod-$6$ content fixed by a single
combinatorial datum: $\{0,1,3\}$ is the \emph{unique} (up to $D_6$) $3$-subset of $C_6=|R^+(A_3)|$ with
all distinct cyclic distances $\{1,2,3\}$, and the three sectors are its $D_6$ orbit
($u\!\equiv\!\{0,1,3\}$, $d\!\equiv\!2{-}\{0,1,3\}$, $e\!\equiv\!2{+}\{0,1,3\}$). The total is fixed by
the carrier, $\sum L=40=|R(D_5)|$, with winding unit $6=|R^+(A_3)|$. Given the residue \emph{set} of a
sector, the full ordered packet is reproduced by two universal rules --- the lightest generation
carries the single $+6$ winding, and the entries are sorted in decreasing length (monotone hierarchy).
The \emph{only} remaining datum is then, per sector, which residue is the first-generation anchor that
receives the winding:
\begin{equation}
  \boxed{\ \text{anchors}\ (a_u,a_d,a_e)=(1,1,2)\ \Rightarrow\
  L_u=(7,3,0),\ L_d=(7,5,2),\ L_e=(8,5,3)\ }\qquad\mE\ (\text{given anchors}) .
\end{equation}
So the open step of a full ``Clebsch walk'' theorem is sharply isolated: it is exactly these three
per-sector anchors. They are \emph{not} free: the three transport cusps are exactly the $D_5$
dual-root combinations, and each sector is tied to its cusp by the right-handed hypercharge,
\begin{equation}
  \boxed{\ \{1,\tfrac23,\tfrac13\}=\{2y_+,\,-2y_-,\,2(y_++y_-)\},\qquad
  \text{lepton}\leftrightarrow 2y_+,\ \text{up}\leftrightarrow -2y_-,\ \text{down}\leftrightarrow 2(y_++y_-)\ }\qquad\mE .
\end{equation}
The canonical-transport-kernel theorem (document~3, Part~II) derives the residue packets --- and hence the anchors ---
as the \emph{word-lengths of the family-connection holonomy} $\operatorname{Hol}_{\mathcal G_{D_4}}(\nabla_F^\star)$
along the $D_4$-invariant geodesic spine of $\mathbb P^1\setminus\mu_4$ with these cusps as boundary
data. The anchors are therefore a genuine $D_5\times A_3$ readout (cusps from the $D_5$ carrier, spine
holonomy from the $A_3$ four-punctured sphere), not an extra parameter. A compact candidate closed form
is the nearest-integer cusp map
\begin{equation}
  \boxed{\ a_f=\operatorname{round}(2\,\mathrm{cusp}_f):\quad (a_u,a_d,a_e)=\operatorname{round}\bigl(\tfrac43,\tfrac23,2\bigr)=(1,1,2)\ }\qquad\mC ,
\end{equation}
which reproduces the matrix exactly. The remaining open step is now precise and purely geometric:
an \emph{independent} closed-form evaluation of the $D_4$-spine holonomy word-lengths from the lattice
data alone (the same $\mathbb P^1\setminus\mu_4$ whose residue pairing is $\mathrm{Cartan}(A_3)$ and
whose corner rotation is the $A_3$ Coxeter element), rather than from the family-connection
theorem (document~3).

\par\smallskip\noindent\textbf{Riemann--Hilbert check, and what it does (and does not) fix.}
Setting up the rigid $\mathbb Z_4$-symmetric rank-$3$ $SU(3)$ local system explicitly --- four puncture
monodromies $M_k=R^kCR^{-k}$ with the order-$3$ cusp class $C\sim\operatorname{diag}(1,\omega,\omega^2)$
($\omega=e^{2\pi i/3}$, the three cusps as eigenvalue-exponents) and the rotation $R$ ($R^4=\mathbb 1$)
--- the closure $M_0M_1M_2M_3=\mathbb 1$ has a solution to machine precision, and the unique $R$ has
eigenvalues $\{i,-1,-i\}$: the $A_3$ Coxeter element. This \emph{independently} confirms the rigidity
of the family connection and the corner-rotation $=$ Coxeter identity. However, the flat
monodromy \emph{alone} does not fix the integer word-lengths $\{0,1,3\}$: those are \emph{parabolic
degrees} (the cusp exponents as parabolic weights, together with the $\mathbb Z_2$ sheet parity), i.e.\
Hodge/metric data on the spine, not topological monodromy. \emph{Status update (the computation has
since been carried out):} the exact residue normal form $A_0^\star$ with
$\operatorname{spec}=\{0,\tfrac13,\tfrac23\}$ is pinned algebraically
(\veri{v115\_anchor\_residue.py}, \veri{v116\_resonance\_uniqueness.py}), the unitarised holonomy is the
exact $24$-element $W(A_3)=S_4$ (\veri{v117\_monodromy\_weyl\_a3.py}), and the parabolic weights are the
$\mu_4$-character degrees of $H^1(\mathbb P^1\setminus\mu_4)$ --- one spectrum in four readouts
(\veri{v137\_qplus\_cohomology.py}); the deck itself is \emph{derived} --- the integer model selects the
sheet-twisted class, no discrete freedom left (\veri{v141\_deck\_selection.py}), and the full M\"obius
$D_4$ of the curve matches the integer model parity by parity (\veri{v146\_moebius\_d4.py}). What
remains of the old residual is the much smaller selector establishment R4$'$ (research contracts):
derive the \emph{one} lift map of the torsion normal --- the pairing values are atom identities
(\veri{v142\_frame\_integrality.py}/\veri{v145\_pairing\_atoms.py}). \mE/\mC
\par\smallskip

\begin{figure}[ht]
\centering
\begin{tikzpicture}[scale=0.92,>=Stealth]
 % --- up sector ---
 \begin{scope}[xshift=0cm]
   \foreach \i in {0,...,5} \coordinate (u\i) at ({90-60*\i}:1.25);
   \draw[gray!45] (u0)--(u1)--(u2)--(u3)--(u4)--(u5)--cycle;
   \foreach \i in {0,...,5} {\fill[gray!25] (u\i) circle (2.3pt); \node[font=\scriptsize,gray!60!black] at ({90-60*\i}:1.62) {\i};}
   \foreach \i in {0,1,3} \fill[orange!80!black] (u\i) circle (3pt);
   \draw[orange!80!black,line width=0.9pt] (u0)--(u1)--(u3)--cycle;
   \node[font=\scriptsize] at (0,-2.0) {up $\equiv\{0,1,3\}$};
 \end{scope}
 % --- down sector ---
 \begin{scope}[xshift=4.3cm]
   \foreach \i in {0,...,5} \coordinate (d\i) at ({90-60*\i}:1.25);
   \draw[gray!45] (d0)--(d1)--(d2)--(d3)--(d4)--(d5)--cycle;
   \foreach \i in {0,...,5} {\fill[gray!25] (d\i) circle (2.3pt); \node[font=\scriptsize,gray!60!black] at ({90-60*\i}:1.62) {\i};}
   \foreach \i in {1,2,5} \fill[teal!70!black] (d\i) circle (3pt);
   \draw[teal!70!black,line width=0.9pt] (d1)--(d2)--(d5)--cycle;
   \node[font=\scriptsize] at (0,-2.0) {down $\equiv 2{-}\{0,1,3\}$};
 \end{scope}
 % --- lepton sector ---
 \begin{scope}[xshift=8.6cm]
   \foreach \i in {0,...,5} \coordinate (e\i) at ({90-60*\i}:1.25);
   \draw[gray!45] (e0)--(e1)--(e2)--(e3)--(e4)--(e5)--cycle;
   \foreach \i in {0,...,5} {\fill[gray!25] (e\i) circle (2.3pt); \node[font=\scriptsize,gray!60!black] at ({90-60*\i}:1.62) {\i};}
   \foreach \i in {2,3,5} \fill[purple!70!black] (e\i) circle (3pt);
   \draw[purple!70!black,line width=0.9pt] (e2)--(e3)--(e5)--cycle;
   \node[font=\scriptsize] at (0,-2.0) {lepton $\equiv 2{+}\{0,1,3\}$};
 \end{scope}
\end{tikzpicture}
\caption{Flavor transport on the hexagon $C_6=Z_6=|R^+(A_3)|$. The unique distinct-distance triple
$\{0,1,3\}$ (cyclic distances $\{1,2,3\}$) and its $D_6$ orbit give the three sectors' residues mod $6$;
the total $\sum L=40=|R(D_5)|$ and the first-generation $+6$ winding come from the carrier.}
\label{fig:hexagon}
\end{figure}

